\section{Results}
\label{sec:results}

\subsection{Main test (Stage~3)}
Table~\ref{tab:main} reports Stage~3 test-only S-phase metrics (10{,}000 traces /
1{,}482 events).
Fixed re-ranking improves F1@$0.5$ from $0.695$ to $0.718$ and reduces e2e P95
from $2.15$\,s to $1.11$\,s ($\approx 48\%$ relative P95 reduction; absolute
values retained).
F1@$0.1$ also rises ($0.462\rightarrow 0.477$), but the larger operational gain
is in suppressing gross wrong-peak tails rather than claiming sub-sample-point
timing refinement.

\begin{table}[t]
\centering
\caption{Stage~3 event-disjoint test-only S-phase results (frozen STEAD PhaseNet).}
\label{tab:main}
\begin{tabular}{lrrrr}
\toprule
Method & F1@0.1 & F1@0.5 & e2e P95 (s) & e2e MAE (s) \\
\midrule
PhaseNet (STEAD) & 0.462 & 0.695 & 2.15 & 1.99 \\
fixed catalog\_rescore & 0.477 & 0.718 & 1.11 & 0.62 \\
learned gate (ablation) & 0.475 & 0.716 & 1.16 & 0.59 \\
oracle PN/fixed selector & 0.477 & 0.720 & 1.10 & 0.62 \\
oracle candidate K=5 & 0.588 & 0.890 & 0.96 & 0.31 \\
\bottomrule
\end{tabular}

\end{table}

Event-level bootstrap of fixed versus PhaseNet (frozen predictions, seed~$42$,
$n{=}2000$) yields $\Delta$F1@$0.5{=}0.0225$ with 95\% CI $[0.019,\,0.026]$ and
$\Delta$P95 with CI entirely below zero
(Table~\ref{tab:boot}; \texttt{bootstrap\_fixed\_vs\_phasenet.json}).
Learned gate versus fixed has $\Delta$F1@$0.5$ CI $[-0.004,\,-0.0004]$ (gate
slightly worse).

\begin{table}[t]
\centering
\caption{Stage~3 event-level paired bootstrap: fixed re-ranking minus PhaseNet.}
\label{tab:boot}
\begin{tabular}{lrrr}
\toprule
Metric & mean $\Delta$ & 95\% CI low & 95\% CI high \\
\midrule
F1@0.5 & 0.0225 & 0.0186 & 0.0265 \\
F1@0.1 & 0.0144 & 0.0113 & 0.0174 \\
e2e P95 (s) & -1.0698 & -1.7000 & -0.6250 \\
e2e MAE (s) & -1.3597 & -1.6208 & -1.1101 \\
wrong-peak rate & -0.0275 & -0.0323 & -0.0227 \\
\bottomrule
\end{tabular}

\end{table}

Figure~\ref{fig:cdf} shows error CDFs; fixed dominates the long tail.
Figure~\ref{fig:f1tol} reports F1 at $0.1$\,s and $0.5$\,s.
Figure~\ref{fig:subset} shows Stage~2 development hard-subset gains (motivational;
not a test claim).
Figures~\ref{fig:caseok}--\ref{fig:casebad} illustrate a wrong-peak repair and a
failure mode from frozen Stage~3 case dumps.

\begin{figure}[t]
\centering
\includegraphics[width=0.72\linewidth]{fig_error_cdf_stage3}
\caption{Stage~3 S-phase absolute-error CDF: PhaseNet vs fixed re-ranking.}
\label{fig:cdf}
\end{figure}

\begin{figure}[t]
\centering
\includegraphics[width=0.72\linewidth]{fig_f1_tolerance_stage3}
\caption{Stage~3 S F1 versus tolerance.}
\label{fig:f1tol}
\end{figure}

\begin{figure}[t]
\centering
\includegraphics[width=0.85\linewidth]{fig_subset_gains_stage2_dev}
\caption{Stage~2 development hard-subset $\Delta$F1@$0.5$ (mixed fixed-eval; not main test).}
\label{fig:subset}
\end{figure}

\begin{figure}[t]
\centering
\includegraphics[width=0.9\linewidth]{fig_case_success_wrongpeak_fix}
\caption{Qualitative success: PhaseNet wrong peak corrected by fixed re-ranking (frozen case dump).}
\label{fig:caseok}
\end{figure}

\begin{figure}[t]
\centering
\includegraphics[width=0.9\linewidth]{fig_case_failure_no_candidate}
\caption{Qualitative failure: no correct candidate in top-$K$ (oracle-candidate gap).}
\label{fig:casebad}
\end{figure}

\subsection{Auxiliary unused-event audit (Stage~4)}
\label{sec:stage4}
On the nine previously unused test events (108 traces), fixed re-ranking moves
S F1@$0.5$ from $0.678$ to $0.690$ and P95 from $1.03$\,s to $0.73$\,s, with the
same direction for ETHZ/SCEDC weights and exact PhaseNet match under
history-unavailable scoring.
However, event-level bootstrap CIs include zero; the audit is underpowered and
is reported only as a directional unused-event check
(\emph{directionally replicated in an underpowered unused-event audit}), not as
independent significance validation.

\subsection{Ablations and ceilings}
Oracle candidate-$K{=}5$ reaches F1@$0.5{=}0.890$, while oracle PN/fixed selector
only $0.720$, indicating that (i)~candidate coverage still limits performance and
(ii)~a binary gate between PhaseNet and fixed has little headroom---consistent
with the learned gate collapsing toward PhaseNet.
Shuffled/biased history controls degrade relative to correct history on
development and Stage~4 audits.
